\documentclass[12pt, a4paper]{article}

\usepackage[margin=2.5cm]{geometry}
\usepackage{amsmath}
\usepackage{amssymb}
\usepackage{graphicx}
\usepackage{enumitem}
\usepackage{hyperref}
\usepackage{ctex}

% 设置段落间距
\setlength{\parskip}{1em}

% 设置文档标题
\title{\textbf{实验报告：基于小波变换的图像压缩}}
\author{} % 空作者
\date{}   % 空日期

\begin{document}

\maketitle % 显示标题

\section*{\centering 一、实验目的}
\begin{enumerate}
    \item 理解小波变换（\textit{Wavelet Transform}）的基本原理及其在图像处理中的应用。
    \item 掌握使用 MATLAB 进行图像基于二维离散小波变换（2D-DWT）的分解与重构方法。
    \item 探究小波变换的能量集中特性，并基于此实现一个基本的图像压缩演示。
    \item 分析不同分解层次对图像压缩效果的影响。
\end{enumerate}

\section*{\centering 二、实验原理与概念介绍}
\subsection*{\centering 1. 图像压缩的基本概念}
图像压缩是指减少表示一幅数字图像所需数据量的技术。其目的是在不（或在可接受范围内）降低图像质量的前提下，去除数据冗余，以便于图像的存储和传输。

图像压缩分为两类：
\begin{itemize}
    \item \textbf{无损压缩（Lossless Compression）}： 压缩和解压缩后的图像与原始图像完全一致，不丢失任何信息。
    \item \textbf{有损压缩（Lossy Compression）}： 用牺牲程序对全部一些信息，解压缩后的图像与原始图像不完全一致，但却能够消除大量冗余，因此压缩率通常远比无损压缩所能达到的要高。
\end{itemize}
本实验所探究的小波压缩是一种典型的有损压缩技术，其基本流程包括三个步骤：\textbf{变换}、\textbf{量化}和\textbf{编码}。

\subsection*{\centering 2. 小波变换（Wavelet Transform）}
\subsubsection*{(1) 为什么需要小波变换？}
传统的傅里叶变换（\textit{Fourier Transform}）是强大的频域分析工具，但它只提供信号的频率信息，而完全丢失了时间（或空间）信息。对于像图像这类非平稳信号（其统计特性随空间位置变化，如边缘），我们不仅想知道有哪些频率分量，更想知道这些频率分量出现在什么位置。

小波变换是一种\textbf{时-频（空-频）}分析方法。它使用一组具有有限时长和振荡特性的小波基函数（\textit{Wavelets}）来表示信号。这些小波基函数可以通过对一个“母小波”进行\textbf{伸缩（Scaling）}和\textbf{平移（Translation）}得到。

伸缩对应\textbf{频率}：对母小波进行拉伸（尺度变大），对应低频信息；进行压缩（尺度变小），对应高频信息。
\textbf{平移对应时间/空间}：将小波函数在信号上滑动，以捕捉不同位置的特征。

因此，小波变换能够同时在空间和频率上定位图像的特征，这使其在图像去噪、边缘检测和压缩领域具有巨大优势。

\subsubsection*{(2) 二维离散小波变换（2D-DWT）}
对于二维的图像，2D-DWT 通常通过可分离的方式实现：首先对图像的每一行进行一维小波变换，然后对结果的每一列进行一维小波变换。

一次 2D-DWT 分解（\textit{Analysis}）会将图像（设为 X）分解为四个大小为原始图像 1/4 的子带：
\begin{itemize}
    \item \textbf{LL 子带（低频分量 / 近似）}：对行和列都进行低通滤波。它保留了图像的主要轮廓和能量，通常是图像的缩小版本。
    \item \textbf{HL 子带（水平细节）}：对行进行高通滤波，对列进行低通滤波。它捕捉了图像的水平方向边缘信息。
    \item \textbf{LH 子带（垂直细节）}：对行进行低通滤波，对列进行高通滤波。它捕捉了图像的垂直方向边缘信息。
    \item \textbf{HH 子带（对角细节）}：对行和列都进行高通滤波。它捕捉了图像的对角线方向边缘和细节。
\end{itemize}

\subsubsection*{(3) 小波分解（Multi-Level Wavelet Decomposition, MRA）}
小波分解允许我们对图像进行多级分解（\textit{Multi-Resolution Analysis, MRA}）。由于图像的大部分能量（即视觉上的主要信息）都集中在低频 LL 子带中，我们可以对这个 LL 子带再次进行上述的四子带分解，得到第二级的 LL2, HL2, LH2, HH2。

这个过程可以重复多次，每分解一次，分辨率就下降一次，形成一个金字塔结构。随着分解层数的增加，LL 子带变得越来越小，而细节系数则被集合包含于图像的“精华为”信息。

\subsection*{\centering 3. 基于小波变换的图像压缩原理}
小波变换之所以能用于压缩，关键在于其优异的能量集中特性：
\begin{itemize}
    \item \textbf{能量集中}： 经过 DWT 后，图像的绝大部分能量集中到少数几个低频系数（LL 子带）中。
    \item \textbf{系数稀疏}： 而三个高频子带（HL, LH, HH）中的系数（代表图像的边缘和纹理）通常幅值很小，大部分接近于零。
\end{itemize}
压缩正是利用了这一特性：
\begin{itemize}
    \item \textbf{变换}： 对图像进行多级 DWT。
    \item \textbf{量化}： 这是实现压缩和引入损失的关键步骤。我们对变换后的系数进行处理：
    \begin{itemize}
        \item \textbf{保留低频}： 对能量集中的低频（LL）系数，我们使用较高的精度来保留它们，因为它们对图像的视觉质量至关重要。例如，对于 2 级分解，我们可以保留 LL2 子带；而 HL2, LH2, HH2 等高频分量的系数直接置为 0，对剩下梯度的系数也使用较低的精度来存储。
    \end{itemize}
    \item \textbf{编码}： 对量化后的系数进行编码。利用霍夫曼编码、算术编码等，进一步减小数据量。
\end{itemize}
本实验通过只提取和显示低频近似分量（\texttt{ca1} 和 \texttt{ca2}），直观地展示了压缩的核心思想：仅用低频分量（占原数据量 1/4 或 1/16）就可以近似重构原始图像。

\subsection*{\centering 4. Bior3.7 小波}
\textbf{Bior3.7} 是 \textbf{Biorthogonal}（双正交）小波族的一种，与 \textbf{Orthogonal}（正交）小波不同，它允许分解小波和重构小波不相同。Bior3.7 小波对于图像的边缘保持非常有效，它常用于图像的压缩和去噪。

\textbf{Bior3.7} 小波具有较好的\textbf{线性和对称性}，避免了图像处理中产生的\textbf{边缘失真}。

使用 Bior3.7 小波进行二维小波分解时，它会生成四个子带（LL, HL, LH, HH）。

\section*{\centering 三、实验环境}
MATLAB

\section*{\centering 四、实验内容与步骤}
\begin{enumerate}
    \item 导入工作区：使用 \texttt{clc}, \texttt{clear all}, \texttt{close all} 清理命令行窗口、工作区变量和关闭所有图像窗口。
    \item 读取图像：使用 \texttt{imread} 读取 \texttt{im.jpg} 图像文件，并使用 \texttt{rgb2gray} 将其转换为灰度图像。
    \item 显示原始图像：使用 \texttt{subplot(221)} 显示原始图像，并在第一个位置 \texttt{(221)} 显示原始灰度图像。
    \item \textbf{一级小波分解}：使用 \texttt{wavedec2} 函数，选择 \texttt{bior3.7} 小波，对图像 \texttt{X} 进行 2 级分解，得到 2 级分解系数向量，并显示分解结果。
    \item \textbf{第 1 级图像分解（即 1 级 DWT）}：
    \begin{itemize}
        \item 使用 \texttt{appcoef2} 函数提取第 1 级低频系数 \texttt{ca1}。
        \item 使用 \texttt{detcoef2} 函数提取第 1 级水平(\texttt{ch1}), 垂直(\texttt{cv1})和对角(\texttt{cd1})高频系数。
        \item 使用 \texttt{wmax} 函数，在 \texttt{subplot(222)} 显示分解结果。
    \end{itemize}
    \item \textbf{显示第 1 级分解结果}： 将重构的四个分量 A1, H1, V1, D1 排成一个矩阵 A1, H1, V1, D1，并在一张图中显示，以便于查看图像分解为低频近似和三个方向的高频细节。
    \item \textbf{第 1 级压缩和解压}：
    \begin{itemize}
        \item 使用 \texttt{appcoef2} 函数提取第 1 级低频系数 \texttt{ca1}。
        \item 然后，将 \texttt{ch1}, \texttt{cv1}, \texttt{cd1} 的系数设置为 0（即为高频分量置零）。
        \item 使用 \texttt{idwt2} 函数，将 \texttt{ca1}, \texttt{ch1}, \texttt{cv1}, \texttt{cd1} 重构为一幅图像。
        \item 在 \texttt{subplot(223)} 中显示 \texttt{ca1} 的图像。
        \item 在 \texttt{subplot(224)} 中显示压缩后的图像。
    \end{itemize}
    \item \textbf{第 2 级压缩和解压}：
    \begin{itemize}
        \item 使用 \texttt{appcoef2} 提取第 2 级低频系数 \texttt{ca2}。
        \item 将所有高频系数置为 0。
        \item 在 \texttt{(224)} 位置显示 \texttt{ca2}。这代表了仅使用 1/16 数据量（因为是 2 级分解）对图像的压缩。
    \end{itemize}
\end{enumerate}

\section*{\centering 五、实验结果与分析}
\subsection*{\centering 实验结果}
实验共生成一个包含四个子图的图像。
\begin{enumerate}
    \item \textbf{子图 (221) - 原始图像}： 显示了原始灰度图像，作为对比基准。
    \item \textbf{子图 (222) - 第 1 级分解后的图像}：
    \begin{itemize}
        \item 左上角 (A1)：第 1 级低频重构，是原始图像的模糊子版。
        \item 右上角 (H1)：第 1 级水平高频重构，清晰地显示了图像中的垂直边缘。
        \item 左下角 (V1)：第 1 级垂直高频重构，清晰地显示了图像中的水平边缘。
        \item 右下角 (D1)：第 1 级对角高频重构，显示了对角线方向的细节。
    \end{itemize}
    \item \textbf{子图 (223) - 第 2 级压缩图像}： 显示第 2 级的低频近似系数 \texttt{ca2}。该图像在视觉上是原始图像的模糊版。原始图像的大小为 256x256，子带 \texttt{ca2} 的大小为 64x64，仅用 1/16 数据量压缩。
    \item \textbf{子图 (224) - 第 2 次压缩图像}： 显示了第 2 级的低频近似系数 \texttt{ca2}。该图像 \texttt{ca2} 更加模糊和缩小，仅保留了最核心的轮廓信息。这模拟了 16:1 的压缩效果（仅保留 LL2）。
\end{enumerate}

\subsection*{\centering 结果分析}
\begin{enumerate}
    \item \textbf{能量集中特性验证}： 对比子图 (222) 中的四个分量，可以清晰地看到，图像的绝大部分视觉信息（主要轮廓、亮度分布）都集中在低频近似（A1）中，而高频分量（H1, V1, D1）则主要包含边缘和纹理信息，且大部分区域接近于黑色（即系数接近 0）。这充分验证了小波变换的能量集中特性。
    \item \textbf{多级分解与压缩}： 对比子图 (223) 和 (224) 可以发现：
    \begin{itemize}
        \item 第 1 级低频 \texttt{ca1} 相当于原始图像的 1/4，但图像已经模糊。
        \item 第 2 级低频 \texttt{ca2} 压缩更严重，图像更模糊，但它占用的数据量更小（仅为原始图像的 1/16），实现了更高的压缩比。
    \end{itemize}
    这表明，在小波压缩中，分解的\textbf{层数}是一个关键参数。层数越多，低频分量（LLn）越小，压缩比越高，但图像质量损失也越大。
    \item \textbf{代码分析}： 实验代码通过 \texttt{appcoef2} 提取低频分量来演示压缩思想。一个完整的图像压缩通常包括：
    \begin{itemize}
        \item 对高频系数 \texttt{ch1}, \texttt{cv1}, \texttt{cd1} 以及更高层的高频系数 (\texttt{ch2}, \texttt{cv2}, \texttt{cd2}...) 进行\textbf{阈值处理和量化}。
        \item 对量化后的所有系数（低频和高频）进行\textbf{熵编码}（如游程编码、霍夫曼编码）。
        \item 在 MATLAB 中，可以执行 \texttt{wavedec2}，\texttt{waverec2} 函数以及所有的 \texttt{dwt} 和 \texttt{idwt} 函数。
    \end{itemize}
    为实现压缩，我们通常会对系数进行进一步处理，如阈值处理和量化，然后进行熵编码。
\end{enumerate}

\section*{\centering 六、实验结论与心得}
\subsection*{\centering 结论：}
\begin{enumerate}
    \item 小波变换能有效地将图像分解为低频（近似）和高频（细节）分量。
    \item 图像的主要能量集中在低频分量中，而高频分量系数稀疏。
    \item 通过仅保留低频分量并对高频分量进行阈值或置零，能有效实现有损图像压缩。
    \item 分解的层数越高，压缩比越高，但图像的保真度越低。
\end{enumerate}

\subsection*{\centering 心得体会：}
本次实验让我对小波变换的“时-频”局部化特性有了更深刻的理解。与傅里叶变换的“全局”特性不同，小波能精准地捕捉图像的局部特征（如边缘），并将其分离到高频子带。这正是它在图像压缩和去噪领域发挥优势的关键。

实验中仅演示了压缩的核心原理，在后续学习中，我将尝试实现一个包含量化和重构的完整压缩/解压缩流程，并引入 PSNR（峰值信噪比）和 CR（压缩比）等客观指标来定量评估压缩效果。

\end{document}